\section{The Rights of Man}

The Rights of Man (see below), expresses the fundamental world view of modernity. There are those who believe it is a debatable question and that the more reasonable choice will prevail. However, there are fundamental reasons why that is not true. Tradition is not at all a political option, unlike various forms of modernity. We can schematize the perspectives as shown in Table \ref{tab:RightsMan}.

\begin{table}[h]\small
\begin{tabular}{ccc}\toprule
 &
\textbf{Tradition} &
\textbf{Modernity}\\\midrule
World view &
Supernatural &
Natural\\
Rationality &
Knowledge of cosmic law &
Instrumental to achieve goals\\
Social distinctions &
Real &
Conventional\\
Natural inequality &
Just &
Unjust\\
Decision making &
Logical &
Psychological\\\bottomrule
\end{tabular}
\label{tab:RightsMan}
\caption{Rights of Man}
\end{table}
As we take the \textbf{French Revolution} as the tipping point separating modernity from tradition, its dogmas can be found in the Declaration on the Rights of Man.

\paragraph{Article 1}
\begin{quotex}
Men are born and remain free and equal in rights. Social distinctions may be founded only upon the general good.

\end{quotex}
Tradition is based on two related principles, which are both denied by article 1.

\begin{itemize}
\item Spiritual authority

A traditional society is based on some spiritual authority which determines what is true and good. Man is born either in sin or under illusion and can only become free after initiation into the rites of the city. The declaration states that man is born free and not subject to anyone else, hence he himself takes the place of God or the higher authority. 
\item Temporal power

The traditional society is based on a hierarchy of increasing power. This is understood as natural or divinely ordained. The declaration, on the contrary, admits social distinctions only insofar as the serve the common good. Thus, such distinctions are considered conventional, not natural. 
\end{itemize}
\paragraph{Article 3}
\begin{quotex}
The principle of all sovereignty resides essentially in the nation. No body nor individual may exercise any authority which does not proceed directly from the nation. 

\end{quotex}
This emphasizes article 1. There is no sovereignty beyond the nation, in particular, there is no spiritual authority higher than the laws of the nation. A fortiori, there is no church or spiritual organization beyond the nation.

\paragraph{Article 4}
\begin{quotex}
Liberty consists in the freedom to do everything which injures no one else; hence the exercise of the natural rights of each man has no limits except those which assure to the other members of the society the enjoyment of the same rights. These limits can only be determined by law.

\end{quotex}
In the traditional view, there is a cosmic law, sometimes called the natural law, which binds the liberty of man. Unlike the mineral, vegetable, and animal kingdoms, which are determined by the natural law, man has the free will to either follow or reject the cosmic law. Article 4 states that man's actions are totally conventional, oblivious to any higher law. A man's actions need not even be rational, which is itself considered an artificial construct limiting his liberty.

\paragraph{Article 6}
\begin{quotex}
Law is the expression of the general will. Every citizen has a right to participate personally, or through his representative, in its foundation. It must be the same for all, whether it protects or punishes. All citizens, being equal in the eyes of the law, are equally eligible to all dignities and to all public positions and occupations, according to their abilities, and without distinction except that of their virtues and talents. 

\end{quotex}
The \textbf{general will} here takes the place of the will of God. In tradition, distinctions are viewed as a special calling or a birth rite. Article 6 regards distinctions as conventional and acceptable insofar as they serve the general will.

\paragraph{Article 10}
\begin{quotex}
No one shall be molested on account of his opinions, including his religious views, provided their manifestation does not disturb the public order established by law. 

\end{quotex}
This article guarantees freedom of religion. Since man is born free and is his own sovereign, his beliefs cannot be limited. However, by article 2, his religious beliefs cannot be used to restrict others, hence the outer effects of such beliefs can be limited by law. In its statement, the American Bill or Rights seems more absolutist, but religious manifestations are still limited by law in the USA.

This articles differs profoundly from the Traditional view, any hint of which will raise all the forces against it. In the traditional view, ``Error has no rights", so religious views outside the spiritual authority have no guarantee, although they may be tolerated in order to maintain public order. This was certainly the case during the era of Christendom. Despite current misconceptions, even the ancient pagan world did not support complete freedom of religions.

\paragraph{Article 11}
\begin{quotex}
The free communication of ideas and opinions is one of the most precious of the rights of man. Every citizen may, accordingly, speak, write, and print with freedom, but shall be responsible for such abuses of this freedom as shall be defined by law. 

\end{quotex}
The same comments to article 10 apply here also. Freedom of speech is considered an advance over the Medieval period. Nevertheless, there is an ``out" in that abuses of such freedom may be restricted by law. Thus we seen in many Western countries in Europe or Canada have laws restricting speech. Interesting, not even truth is the standard, but rather public order.

So the problem that was supposed to have been solved in the Enlightenment still exists, to wit, who is the ultimate arbiter of speech. Since in the modern state, not even rationality can restrict man's liberty, so similarity conformance of speech to truth is not of ultimate import, even though the very goal of the Enlightenment is truth!

The traditional society regarded its foundation, laws, and mores as infallible. The modern state does not, nor can it do so. Actually, it is irrelevant, since the general will is the sole standard for law and is not subject to any higher truth, whether of a scientific, philosophic, religious, or metaphysical nature. This is critical to understand for anyone engaged in public debate in the modern world.



\flrightit{Posted on 2011-10-26 by Cologero }
